% ============================================================
%  SOAL MATEMATIKA  --  SSU-ITB 2024
%  Tipe A (1-20) + Tipe B (21-39) + Soal Baru (40)
%  Paket 1  |  PG Biasa 16 (40%) | PG Majemuk 16 (40%) | Isian 8 (20%)
% ============================================================

\documentclass[11pt]{article}

\usepackage[utf8]{inputenc}
\usepackage[T1]{fontenc}
\usepackage[a4paper,margin=2.5cm,top=2.8cm]{geometry}
\usepackage{amsmath,amssymb}
\usepackage{graphicx}
\usepackage{enumitem}
\usepackage{multicol}
\usepackage{xcolor}
\usepackage[most]{tcolorbox}
\usepackage{fancyhdr}
\usepackage{booktabs}
\usepackage{icomma}
\usepackage{array}

\pagestyle{fancy}\fancyhf{}
\lhead{\small SSU-ITB 2024}
\rhead{\small Paket 1}
\cfoot{\small\thepage}
\renewcommand{\headrulewidth}{0.4pt}

\newtcolorbox{soal}[1]{
  breakable,
  colback=white, colframe=black, colbacktitle=black,
  coltitle=white, fonttitle=\bfseries\sffamily,
  title=Nomor #1,
  boxrule=0.8pt, arc=0pt,
  left=3mm, right=3mm, top=2mm, bottom=2mm}

\newenvironment{pil}{%
  \begin{enumerate}[label=(\Alph*),leftmargin=*,itemsep=1pt,topsep=3pt]}%
  {\end{enumerate}}
\newenvironment{pilB}{%
  \par\smallskip\begin{multicols}{2}
  \begin{enumerate}[label=(\Alph*),leftmargin=*,itemsep=1pt,topsep=3pt]}%
  {\end{enumerate}\end{multicols}}

\newcommand{\gambar}[2][0.52\textwidth]{%
  \begin{center}\includegraphics[width=#1]{#2}\end{center}}

\linespread{1.05}

% ===================================================================
\begin{document}
\thispagestyle{empty}

\begin{center}
  \fbox{\parbox{0.45\textwidth}{\centering\bfseries\large LEMBAR SOAL}}
\end{center}
\vspace{0.5cm}

\begin{center}
  {\LARGE\bfseries MATEMATIKA}
\end{center}
\vspace{0.5cm}

\begin{center}
  \begin{tabular}{@{\hspace{2cm}}ll@{\hspace{0.4cm}:@{\hspace{0.4cm}}}l}
    Jenjang      & & SMA / Sederajat               \\[0.3em]
    Hari/Tanggal & & \underline{\hspace{5cm}}      \\[0.3em]
    Waktu        & & \underline{\hspace{5cm}}      \\[0.3em]
    Paket Soal   & & 1                             \\
  \end{tabular}
\end{center}

\vspace{0.6cm}
\noindent\rule{\linewidth}{0.6pt}
\vspace{0.3cm}

\noindent\textbf{\underline{PETUNJUK UMUM}}
\begin{enumerate}[leftmargin=*,itemsep=4pt,topsep=4pt]
  \item Anda \textbf{tidak} diperbolehkan menggunakan sumber-sumber eksternal,
        termasuk internet, \textit{large language model} seperti ChatGPT, Claude,
        LLaMA, Gemini, Meta AI, dan sebagainya. Anda sangat dianjurkan untuk
        mencoba mengerjakan sendiri terlebih dahulu karena evaluasi ini dibuat
        untuk \textbf{mengukur} tingkat kepahaman; nilai $<70$ mengindikasikan
        \textbf{tidak lulus}, dan jika sudah melewati \textit{deadline} maka
        dianggap tidak mengerjakan yang berarti \textbf{tidak lulus}.

  \item Terdapat 2 tahap pengerjaan, yaitu tahap \textbf{Try Out} dan tahap
        \textbf{Review}. Try Out bersifat \textit{open book}; Anda diperbolehkan
        membuka buku apapun selama bersifat \textit{offline} dengan rentang waktu
        2 jam pada tanggal \underline{\hspace{3cm}}. Setelah itu memasuki
        tahap Review dengan \textit{schedule} rentang tanggal
        \underline{\hspace{3cm}}--\underline{\hspace{3cm}}, di mana Anda
        diharapkan mengerjakan ulang dengan disertakan penjelasan/pembelaan
        mengapa revisi jawaban adalah pilihan tersebut, dan tetap mencantumkan
        penjelasan pada soal yang walaupun pada penilaian sudah benar.

  \item Anda dianjurkan untuk mencantumkan caranya walaupun tidak termasuk
        penilaian, mengingat sifatnya adalah sebagai bahan pembelajaran.
        Mohon bertanggung jawab atas pekerjaan Anda sendiri.

  \item Bacalah setiap soal dengan teliti. Terdapat tiga jenis soal:
        \begin{enumerate}[label=(\arabic*),leftmargin=2em,topsep=2pt,itemsep=2pt]
          \item \textbf{Pilihan Ganda Biasa} --- 5 opsi (A--E), pilih
                1 jawaban paling tepat;
          \item \textbf{Pilihan Ganda Majemuk Kompleks} --- tabel
                \textbf{Benar}/\textbf{Salah}, tentukan kebenaran tiap
                pernyataan secara mandiri; dan
          \item \textbf{Isian Singkat} --- tulis nilai/ekspresi pada
                kolom yang tersedia.
        \end{enumerate}

  \item Soal berjumlah \textbf{40 butir} soal dengan bobot asli per soal adalah
        \textbf{2{,}5} untuk mendapatkan nilai sempurna sebesar \textbf{100}.

  \item Silakan bertanya apabila terdapat soal yang ambigu, rancu, atau kurang
        spesifik selama itu tidak menjurus kepada jawaban soal tersebut.

  \item Isilah detail informasi berikut.\\[0.3em]
        \textbf{Nama}\hspace{0.45cm};\enspace\underline{\hspace{8cm}}
\end{enumerate}

\vspace{0.6cm}
\noindent\textit{Dengan menandatangani kolom di bawah ini, saya menyatakan telah membaca,
memahami, dan menyetujui seluruh peraturan yang tercantum di atas.}

\vspace{0.6cm}
\noindent\fbox{\parbox{5cm}{\vspace{2.2cm}\centering\small Tanda Tangan Peserta\vspace{0.3cm}}}

\clearpage

% ===================================================================
%  TIPE A  (nomor 1--20)
% ===================================================================

%---------- 1: PG Biasa ----------
\begin{soal}{1}
Jika titik $P(a,b)$ berada pada kurva $y=x^3+5x-7$ di kuadran pertama
sehingga gradien garis singgung kurva tersebut di $P$ adalah 17,
maka nilai $a+b$ adalah \ldots
\begin{pilB}
  \item $9$ \item $11$ \item $13$ \item $15$ \item $17$
\end{pilB}
\end{soal}

%---------- 2: Isian ----------
\begin{soal}{2}
Perhatikan gambar berikut.
\gambar[0.46\textwidth]{images/1/1-17.png}
\begin{enumerate}[label=(\alph*),leftmargin=*,itemsep=6pt,topsep=4pt]
  \item Jika $L(a)$ menyatakan luas persegi panjang di atas,
        maka $L(8)$ adalah \ldots
  \item Jika $c$ memenuhi $L'(c)=0$, maka nilai $2\sqrt{2}\,c$ adalah \ldots
  \item Luas persegi panjang maksimum yang mungkin termuat pada
        daerah setengah lingkaran di atas adalah \ldots
\end{enumerate}
\end{soal}

%---------- 3: PG Biasa ----------
\begin{soal}{3}
Dua bola diambil sekaligus dari suatu kotak yang berisi 5 bola biru
dan 6 bola kuning. Jika peluang terambil paling banyak 1 bola kuning
adalah $\dfrac{A}{11\times10}$, maka nilai $A$ adalah \ldots
\begin{pilB}
  \item $40$ \item $60$ \item $80$ \item $90$ \item $100$
\end{pilB}
\end{soal}

%---------- 4: PG Majemuk Kompleks ----------
\begin{soal}{4}
Diberikan empat pernyataan tentang hiperbola $y^2-x^2=1$.
Tentukan kebenaran setiap pernyataan berikut!

\par\smallskip
{\renewcommand{\arraystretch}{1.4}%
\begin{tabular}{@{}p{0.76\linewidth}cc@{}}
\hline
\textbf{Pernyataan} & \textbf{Benar} & \textbf{Salah}\\
\hline
(1)~Hiperbola $y^2-x^2=1$ memotong sumbu $y$ di dua titik.
  & $\bigcirc$ & $\bigcirc$\\
(2)~Hiperbola $y^2-x^2=1$ merupakan grafik suatu fungsi.
  & $\bigcirc$ & $\bigcirc$\\
(3)~Hiperbola $y^2-x^2=1$ tidak memotong sumbu $x$.
  & $\bigcirc$ & $\bigcirc$\\
(4)~Hiperbola $y^2-x^2=1$ tidak memiliki asimtot miring.
  & $\bigcirc$ & $\bigcirc$\\
\hline
\end{tabular}}
\end{soal}

%---------- 5: PG Biasa ----------
\begin{soal}{5}
Jika $\cos(2x)=-\dfrac{7}{25}$ dengan $0<x<\pi$,
maka $25\sin(x)$ adalah \ldots
\begin{pilB}
  \item $12$ \item $15$ \item $16$ \item $20$ \item $24$
\end{pilB}
\end{soal}

%---------- 6: PG Biasa ----------
\begin{soal}{6}
Diberikan titik $A(1,-3,3)$ dan titik $B(-5,-2,3)$. Jika
$P(p_1,p_2,p_3)$ adalah suatu titik sehingga
$\overrightarrow{BP}=-4\overrightarrow{PA}$,
maka nilai $-3p_2$ adalah \ldots
\begin{pilB}
  \item $-15$ \item $-10$ \item $5$ \item $10$ \item $15$
\end{pilB}
\end{soal}

%---------- 7: PG Majemuk Kompleks ----------
\begin{soal}{7}
Diketahui suku ke-10 dan suku ke-20 suatu barisan aritmetika,
berturut-turut, adalah 84 dan 174.
Tentukan kebenaran setiap pernyataan berikut!

\par\smallskip
{\renewcommand{\arraystretch}{1.4}%
\begin{tabular}{@{}p{0.76\linewidth}cc@{}}
\hline
\textbf{Pernyataan} & \textbf{Benar} & \textbf{Salah}\\
\hline
(1)~Beda $(d)$ barisan tersebut adalah 18.
  & $\bigcirc$ & $\bigcirc$\\
(2)~Suku pertama barisan tersebut adalah 3.
  & $\bigcirc$ & $\bigcirc$\\
(3)~Suku ke-15 barisan tersebut adalah 129.
  & $\bigcirc$ & $\bigcirc$\\
(4)~Jumlah 15 suku pertama $S_{15}=990$.
  & $\bigcirc$ & $\bigcirc$\\
\hline
\end{tabular}}
\end{soal}

%---------- 8: PG Majemuk Kompleks ----------
\begin{soal}{8}
Perhatikan fungsi $f(x)=\sqrt{8-|3x-7|}$.
Tentukan kebenaran setiap pernyataan berikut!

\par\smallskip
{\renewcommand{\arraystretch}{1.4}%
\begin{tabular}{@{}p{0.76\linewidth}cc@{}}
\hline
\textbf{Pernyataan} & \textbf{Benar} & \textbf{Salah}\\
\hline
(1)~$x=0$ termasuk dalam domain fungsi $f$.
  & $\bigcirc$ & $\bigcirc$\\
(2)~$x=5$ termasuk dalam domain fungsi $f$.
  & $\bigcirc$ & $\bigcirc$\\
(3)~$x=-1$ termasuk dalam domain fungsi $f$.
  & $\bigcirc$ & $\bigcirc$\\
(4)~Domain fungsi $f$ adalah $\left[-\tfrac{1}{3},\,5\right]$.
  & $\bigcirc$ & $\bigcirc$\\
\hline
\end{tabular}}
\end{soal}

%---------- 9: PG Majemuk Kompleks ----------
\begin{soal}{9}
Diberikan $f(x)={}^{13}\!\log(x)$.
Tentukan kebenaran setiap pernyataan berikut!

\par\smallskip
{\renewcommand{\arraystretch}{1.4}%
\begin{tabular}{@{}p{0.76\linewidth}cc@{}}
\hline
\textbf{Pernyataan} & \textbf{Benar} & \textbf{Salah}\\
\hline
(1)~$f(5x)+f\!\left(\dfrac{5}{x}\right)={}^{13}\!\log 25$.
  & $\bigcirc$ & $\bigcirc$\\[6pt]
(2)~$f(5x)+f\!\left(\dfrac{5}{x}\right)=2\cdot{}^{13}\!\log 5$.
  & $\bigcirc$ & $\bigcirc$\\[6pt]
(3)~$f(5x)+f\!\left(\dfrac{5}{x}\right)={}^{13}\!\log 10$.
  & $\bigcirc$ & $\bigcirc$\\[6pt]
(4)~$f(169)=2$.
  & $\bigcirc$ & $\bigcirc$\\
\hline
\end{tabular}}
\end{soal}

%---------- 10: PG Majemuk Kompleks ----------
\begin{soal}{10}
Data 90 dari 100 mahasiswa baru suatu perguruan tinggi disajikan
pada tabel berikut.
\begin{center}
\begin{tabular}{l|ccc}
\hline
Jalur   & Kota A & Kota B & Selain A dan B \\
\hline
SNMPTN  & 10 & 12 &  8 \\
UTBK    &  9 &  4 & 17 \\
Mandiri & 11 & 14 &  5 \\
\hline
\end{tabular}
\end{center}
Adapun 10 mahasiswa sisanya bukan dari Kota A dan bukan jalur SNMPTN.
Misalkan $X$ adalah mahasiswa jalur UTBK yang dipilih acak dari 100
mahasiswa. Tentukan kebenaran setiap pernyataan berikut!

\par\smallskip
{\renewcommand{\arraystretch}{1.4}%
\begin{tabular}{@{}p{0.76\linewidth}cc@{}}
\hline
\textbf{Pernyataan} & \textbf{Benar} & \textbf{Salah}\\
\hline
(1)~Dari 90 data, jumlah mahasiswa UTBK asal Kota B adalah 4.
  & $\bigcirc$ & $\bigcirc$\\
(2)~Dari 90 data, setiap jalur memiliki tepat 30 mahasiswa.
  & $\bigcirc$ & $\bigcirc$\\
(3)~Peluang $X$ berasal dari Kota B dapat dipastikan $<\tfrac{1}{12}$.
  & $\bigcirc$ & $\bigcirc$\\
(4)~Dari 90 data, jumlah dari Kota A sama dengan jumlah dari Kota B.
  & $\bigcirc$ & $\bigcirc$\\
\hline
\end{tabular}}
\end{soal}

%---------- 11: PG Majemuk Kompleks ----------
\begin{soal}{11}
Diketahui suatu parabola grafiknya melalui $(-4,0)$ dan $(4,0)$,
serta titik maksimumnya terletak pada garis $y=23x+48$.
Tentukan kebenaran setiap pernyataan berikut!

\par\smallskip
{\renewcommand{\arraystretch}{1.4}%
\begin{tabular}{@{}p{0.76\linewidth}cc@{}}
\hline
\textbf{Pernyataan} & \textbf{Benar} & \textbf{Salah}\\
\hline
(1)~Koefisien $x^2$ pada persamaan parabola bernilai negatif.
  & $\bigcirc$ & $\bigcirc$\\
(2)~Persamaan parabola tersebut adalah $y=-3x^2+48$.
  & $\bigcirc$ & $\bigcirc$\\
(3)~Titik puncak parabola berada di $(0,24)$.
  & $\bigcirc$ & $\bigcirc$\\
(4)~Parabola memotong sumbu $x$ di $(-4,0)$ dan $(4,0)$.
  & $\bigcirc$ & $\bigcirc$\\
\hline
\end{tabular}}
\end{soal}

%---------- 12: PG Biasa ----------
\begin{soal}{12}
Tentukan nilai
\[
  \lim_{x\to 3}\frac{\sin(21-7x)}{x-3}.
\]
\begin{pilB}
  \item $-21$ \item $-7$ \item $0$ \item $7$ \item $21$
\end{pilB}
\end{soal}

%---------- 13: PG Biasa ----------
\begin{soal}{13}
Jika
\[
  \int_{-7}^{-3}f(x)\,dx=12,
\]
maka nilai
\[
  \int_{3}^{5}f(3-2x)\,dx
\]
adalah \ldots
\begin{pilB}
  \item $2$ \item $4$ \item $6$ \item $8$ \item $12$
\end{pilB}
\end{soal}

%---------- 14: PG Majemuk Kompleks ----------
\begin{soal}{14}
Perhatikan grafik fungsi $f$ berikut.
\gambar[0.5\textwidth]{images/1/1-18.png}
Diketahui $a=6$. Tentukan kebenaran setiap pernyataan berikut!

\par\smallskip
{\renewcommand{\arraystretch}{1.4}%
\begin{tabular}{@{}p{0.76\linewidth}cc@{}}
\hline
\textbf{Pernyataan} & \textbf{Benar} & \textbf{Salah}\\
\hline
(1)~$f(6)=36$.
  & $\bigcirc$ & $\bigcirc$\\
(2)~$f'(6)=0$, karena $a=6$ adalah titik maksimum fungsi $f$.
  & $\bigcirc$ & $\bigcirc$\\
(3)~$\displaystyle\lim_{x\to 6}\dfrac{f(x)-6x}{x-6}=f'(6)-6$.
  & $\bigcirc$ & $\bigcirc$\\
(4)~$\displaystyle\lim_{x\to 6}\dfrac{f(x)-6x}{x-6}=6$.
  & $\bigcirc$ & $\bigcirc$\\
\hline
\end{tabular}}
\end{soal}

%---------- 15: PG Majemuk Kompleks ----------
\begin{soal}{15}
Diketahui $x^3-4x^2+9x+a=(x-1)\,g(x)$ dengan $a$ konstanta real.
Tentukan kebenaran setiap pernyataan berikut!

\par\smallskip
{\renewcommand{\arraystretch}{1.4}%
\begin{tabular}{@{}p{0.76\linewidth}cc@{}}
\hline
\textbf{Pernyataan} & \textbf{Benar} & \textbf{Salah}\\
\hline
(1)~Nilai $a=-6$.
  & $\bigcirc$ & $\bigcirc$\\
(2)~Nilai $g(1)=4$.
  & $\bigcirc$ & $\bigcirc$\\
(3)~Polinom $x^3-4x^2+9x-6$ habis dibagi $(x-1)$.
  & $\bigcirc$ & $\bigcirc$\\
(4)~Nilai $g(0)=-6$.
  & $\bigcirc$ & $\bigcirc$\\
\hline
\end{tabular}}
\end{soal}

%---------- 16: PG Majemuk Kompleks ----------
\begin{soal}{16}
Diketahui fungsi $y=6\cos\!\left(\dfrac{\pi t}{3}\right)+9$.
Tentukan kebenaran setiap pernyataan berikut!

\par\smallskip
{\renewcommand{\arraystretch}{1.4}%
\begin{tabular}{@{}p{0.76\linewidth}cc@{}}
\hline
\textbf{Pernyataan} & \textbf{Benar} & \textbf{Salah}\\
\hline
(1)~Periode fungsi tersebut adalah 6.
  & $\bigcirc$ & $\bigcirc$\\
(2)~Amplitudo fungsi tersebut adalah 9.
  & $\bigcirc$ & $\bigcirc$\\
(3)~Nilai minimum fungsi tersebut adalah 3.
  & $\bigcirc$ & $\bigcirc$\\
(4)~Jika $P$ menyatakan periode dan $A$ nilai minimum, maka $10P+A=63$.
  & $\bigcirc$ & $\bigcirc$\\
\hline
\end{tabular}}
\end{soal}

%---------- 17: PG Biasa ----------
\begin{soal}{17}
Jika $f(x)=(a+1)x^3+x^2+12$ dan $f'(1)=32$,
maka nilai $a$ adalah \ldots
\begin{pilB}
  \item $5$ \item $7$ \item $8$ \item $9$ \item $11$
\end{pilB}
\end{soal}

%---------- 18: PG Majemuk Kompleks ----------
\begin{soal}{18}
Diketahui $f(x)=ax^2+b$,\;
$\displaystyle\int_0^1 f(x)\,dx=6$,\;
dan $\displaystyle\int_1^2 f(x)\,dx=12$.
Tentukan kebenaran setiap pernyataan berikut!

\par\smallskip
{\renewcommand{\arraystretch}{1.4}%
\begin{tabular}{@{}p{0.76\linewidth}cc@{}}
\hline
\textbf{Pernyataan} & \textbf{Benar} & \textbf{Salah}\\
\hline
(1)~Nilai $a=3$.
  & $\bigcirc$ & $\bigcirc$\\
(2)~Nilai $b=5$.
  & $\bigcirc$ & $\bigcirc$\\
(3)~$f(1)=8$.
  & $\bigcirc$ & $\bigcirc$\\
(4)~$\displaystyle\int_0^2 f(x)\,dx=24$.
  & $\bigcirc$ & $\bigcirc$\\
\hline
\end{tabular}}
\end{soal}

%---------- 19: PG Majemuk Kompleks ----------
\begin{soal}{19}
Diberikan transformasi $T$ dengan matriks penyajian $M$ ($2\times2$),
yaitu $T(X)=MX$ untuk tiap $X=\begin{bmatrix}x\\y\end{bmatrix}$.
Diketahui
\[
  T\!\left(\begin{bmatrix}3\\-2\end{bmatrix}\right)=\begin{bmatrix}1\\1\end{bmatrix}
  \quad\text{dan}\quad
  T\!\left(\begin{bmatrix}-1\\1\end{bmatrix}\right)=\begin{bmatrix}4\\2\end{bmatrix}.
\]
Jika $T(X)=\begin{bmatrix}9\\5\end{bmatrix}$, tentukan kebenaran
pernyataan-pernyataan berikut tentang $X$!

\par\smallskip
{\renewcommand{\arraystretch}{1.4}%
\begin{tabular}{@{}p{0.76\linewidth}cc@{}}
\hline
\textbf{Pernyataan} & \textbf{Benar} & \textbf{Salah}\\
\hline
(1)~Komponen pertama dari $X$ adalah 1.
  & $\bigcirc$ & $\bigcirc$\\
(2)~Komponen kedua dari $X$ adalah 1.
  & $\bigcirc$ & $\bigcirc$\\
(3)~$X=1\cdot\begin{bmatrix}3\\-2\end{bmatrix}+2\cdot\begin{bmatrix}-1\\1\end{bmatrix}$.
  & $\bigcirc$ & $\bigcirc$\\
(4)~Norma vektor $X$ adalah $|X|=\sqrt{2}$.
  & $\bigcirc$ & $\bigcirc$\\
\hline
\end{tabular}}
\end{soal}

%---------- 20: PG Majemuk Kompleks ----------
\begin{soal}{20}
Diberikan limas persegi $T.ABCD$ seperti pada gambar berikut.
\gambar[0.52\textwidth]{images/1/1-19.png}
Diketahui $|AB|=2\sqrt{2}$ dan tinggi limas adalah 4. Titik $P$
adalah titik tengah $\overline{AT}$ dan titik $Q$ pada $\overline{CT}$
dengan $|CT|=4|CQ|$. Tentukan kebenaran setiap pernyataan berikut!

\par\smallskip
{\renewcommand{\arraystretch}{1.4}%
\begin{tabular}{@{}p{0.76\linewidth}cc@{}}
\hline
\textbf{Pernyataan} & \textbf{Benar} & \textbf{Salah}\\
\hline
(1)~Titik $P$ berada tepat di tengah ruas garis $\overline{AT}$.
  & $\bigcirc$ & $\bigcirc$\\
(2)~$|CQ|=\tfrac{1}{4}|CT|$, sehingga $Q$ seperempat perjalanan dari $C$ ke $T$.
  & $\bigcirc$ & $\bigcirc$\\
(3)~Garis $\overline{PQ}$ sejajar dengan bidang alas $ABCD$.
  & $\bigcirc$ & $\bigcirc$\\
(4)~Tangen sudut antara $\overline{AC}$ dan $\overline{PQ}$ adalah $\tfrac{2}{5}$.
  & $\bigcirc$ & $\bigcirc$\\
\hline
\end{tabular}}
\end{soal}

% ===================================================================
%  TIPE B  (nomor 21--39)
% ===================================================================

%---------- 21: PG Majemuk Kompleks ----------
\begin{soal}{21}
Diberikan garis $y=16-4x$ dan kurva $y=16-x^2$. Titik potong keduanya
adalah $(0,a)$ dan $(b,c)$, serta $L$ adalah luas daerah tertutup
yang dibatasi kurva tersebut.
Tentukan kebenaran setiap pernyataan berikut!

\par\smallskip
{\renewcommand{\arraystretch}{1.4}%
\begin{tabular}{@{}p{0.76\linewidth}cc@{}}
\hline
\textbf{Pernyataan} & \textbf{Benar} & \textbf{Salah}\\
\hline
(1)~Koordinat $y$ titik potong pertama (di $x=0$) adalah $a=16$.
  & $\bigcirc$ & $\bigcirc$\\
(2)~Koordinat $x$ titik potong kedua adalah $b=4$.
  & $\bigcirc$ & $\bigcirc$\\
(3)~Nilai $6L=48$.
  & $\bigcirc$ & $\bigcirc$\\
(4)~Nilai $a+b=20$.
  & $\bigcirc$ & $\bigcirc$\\
\hline
\end{tabular}}
\end{soal}

%---------- 22: Isian ----------
\begin{soal}{22}
Diberikan balok $ABCD.EFGH$ dengan $|AB|=3$ dan $|BC|=|AE|=7$.
\gambar[0.5\textwidth]{images/1/1-20.png}
Misalkan $P$ adalah titik pada $BC$ dan $Q$ pada $CG$ sehingga
$|BP|=|GQ|=1$. Kosinus sudut antara garis $\overline{AP}$ dan
garis $\overline{HQ}$ adalah \ldots
\end{soal}

%---------- 23: PG Majemuk Kompleks ----------
\begin{soal}{23}
Misalkan $s$ adalah simpangan baku dari data terurut enam bilangan:
$4,\ 6-a,\ 6,\ 8,\ 8+a,\ 10$.
Tentukan kebenaran setiap pernyataan berikut!

\par\smallskip
{\renewcommand{\arraystretch}{1.4}%
\begin{tabular}{@{}p{0.76\linewidth}cc@{}}
\hline
\textbf{Pernyataan} & \textbf{Benar} & \textbf{Salah}\\
\hline
(1)~Rata-rata data tersebut adalah 7 untuk setiap nilai $a$.
  & $\bigcirc$ & $\bigcirc$\\
(2)~Nilai varians $s^2$ merupakan fungsi tidak turun terhadap $a\geq0$.
  & $\bigcirc$ & $\bigcirc$\\
(3)~Untuk $a=0$, nilai $s^2<6$.
  & $\bigcirc$ & $\bigcirc$\\
(4)~Untuk setiap $a$ yang memenuhi keterurutan data, selalu berlaku $s^2\leq6$.
  & $\bigcirc$ & $\bigcirc$\\
\hline
\end{tabular}}
\end{soal}

%---------- 24: PG Biasa ----------
\begin{soal}{24}
Jika $r_{1,2}=6\pm\sqrt{2}$ adalah akar-akar persamaan kuadrat
$x^2-Px+Q=0$, maka $P+Q$ adalah \ldots
\begin{pilB}
  \item $34$ \item $38$ \item $42$ \item $46$ \item $50$
\end{pilB}
\end{soal}

%---------- 25: PG Biasa ----------
\begin{soal}{25}
Jika $y=x^2-16x+61$ dengan $x>8$, maka $x=C+\sqrt{y+D}$ dengan
$C$ dan $D$ merupakan dua konstanta bilangan real. Nilai $10C+D$
adalah \ldots
\begin{pilB}
  \item $75$ \item $80$ \item $83$ \item $86$ \item $90$
\end{pilB}
\end{soal}

%---------- 26: PG Biasa ----------
\begin{soal}{26}
Jika $\cos(2x)=-\dfrac{34}{64}$ dengan $0<x<\pi$,
maka $64\sin(x)$ adalah \ldots
\begin{pilB}
  \item $28$ \item $42$ \item $49$ \item $56$ \item $64$
\end{pilB}
\end{soal}

%---------- 27: Isian ----------
\begin{soal}{27}
Diketahui $AB$ adalah diameter suatu lingkaran dengan $A=(10,16)$
dan $B=(-2,0)$. Tentukan jari-jari lingkaran dan titik pusat
lingkaran tersebut.
\end{soal}

%---------- 28: PG Majemuk Kompleks ----------
\begin{soal}{28}
Nilai fungsi $f(x)$, $g(x)$, dan turunannya di $x=0$ dan $x=1$
diberikan pada tabel berikut.
\begin{center}
\begin{tabular}{c|cccc}
$x$ & $f(x)$ & $g(x)$ & $f'(x)$ & $g'(x)$ \\
\hline
$0$ & $2$ & $1$ & $3$ & $\tfrac{2}{3}$ \\
$1$ & $3$ & $-5$ & $-\tfrac{1}{4}$ & $-1$ \\
\end{tabular}
\end{center}
Misalkan $h(x)=\dfrac{2f(x)}{g(x)+1}$.
Tentukan kebenaran setiap pernyataan berikut!

\par\smallskip
{\renewcommand{\arraystretch}{1.4}%
\begin{tabular}{@{}p{0.76\linewidth}cc@{}}
\hline
\textbf{Pernyataan} & \textbf{Benar} & \textbf{Salah}\\
\hline
(1)~$h(0)=2$.
  & $\bigcirc$ & $\bigcirc$\\
(2)~Penghitungan $h'(0)$ menggunakan aturan hasil bagi.
  & $\bigcirc$ & $\bigcirc$\\
(3)~$h'(0)=\dfrac{14}{3}$.
  & $\bigcirc$ & $\bigcirc$\\
(4)~$h'(0)=\dfrac{7}{3}$.
  & $\bigcirc$ & $\bigcirc$\\
\hline
\end{tabular}}
\end{soal}

%---------- 29: Isian ----------
\begin{soal}{29}
Tabel berikut menyajikan nilai fungsi $f$ dan turunannya di
beberapa titik.
\begin{center}
\begin{tabular}{c|cccc}
$x$ & $0$ & $4$ & $5$ & $23$ \\
\hline
$f(x)$ & $23$ & $0$ & $7$ & $0$ \\
$f'(x)$ & $7$ & $8$ & $3$ & $6$ \\
\end{tabular}
\end{center}
Tentukan nilai
\[
  \lim_{x\to 4}\frac{f(5x+3)}{x-4}.
\]
\end{soal}

%---------- 30: PG Majemuk Kompleks ----------
\begin{soal}{30}
Perhatikan limit berikut.
\[
  \lim_{x\to 9}\frac{x+5\sqrt{x}-36}{\sqrt{x}-4}
\]
Tentukan kebenaran setiap pernyataan berikut!

\par\smallskip
{\renewcommand{\arraystretch}{1.4}%
\begin{tabular}{@{}p{0.76\linewidth}cc@{}}
\hline
\textbf{Pernyataan} & \textbf{Benar} & \textbf{Salah}\\
\hline
(1)~Saat $x=9$, penyebut $\sqrt{x}-4=-1\neq0$.
  & $\bigcirc$ & $\bigcirc$\\
(2)~Nilai limit dapat dihitung dengan substitusi langsung $x=9$.
  & $\bigcirc$ & $\bigcirc$\\
(3)~Nilai limit tersebut adalah $-12$.
  & $\bigcirc$ & $\bigcirc$\\
(4)~Nilai limit tersebut adalah $12$.
  & $\bigcirc$ & $\bigcirc$\\
\hline
\end{tabular}}
\end{soal}

%---------- 31: PG Biasa ----------
\begin{soal}{31}
Jika
\[
  \int_5^8 f(x)\,dx=3
  \qquad\text{dan}\qquad
  \int_8^5 g(x)\,dx=5,
\]
maka nilai
\[
  \int_5^8\bigl[3f(x)-2g(x)\bigr]\,dx
\]
adalah \ldots
\begin{pilB}
  \item $1$ \item $9$ \item $12$ \item $19$ \item $24$
\end{pilB}
\end{soal}

%---------- 32: PG Biasa ----------
\begin{soal}{32}
Jika
\[
  x^6-5x^3+6x^2+ax+b=(x^2-1)\,g(x),
\]
maka $10a+b$ adalah \ldots
\begin{pilB}
  \item $35$ \item $39$ \item $43$ \item $47$ \item $51$
\end{pilB}
\end{soal}

%---------- 33: Isian ----------
\begin{soal}{33}
Dari data penerimaan mahasiswa ITB didapatkan informasi berikut.
\begin{enumerate}[leftmargin=*,itemsep=2pt,topsep=3pt]
  \item Sebanyak 46\% dari seluruh mahasiswa pernah mengunjungi
        Kampus Ganesha.
  \item Sebanyak 31\% dari seluruh mahasiswa pernah mengunjungi
        Kampus Jatinangor.
  \item Sebanyak 39\% dari seluruh mahasiswa pernah mengunjungi
        Kampus Ganesha tetapi belum pernah mengunjungi Kampus
        Jatinangor.
\end{enumerate}
Jika dipilih satu mahasiswa secara acak, peluang mahasiswa tersebut
\textbf{belum pernah} mengunjungi Kampus Ganesha maupun Kampus
Jatinangor adalah \ldots
\end{soal}

%---------- 34: PG Biasa ----------
\begin{soal}{34}
Misalkan $x=60\sin t$ dan $y=10\cos t$. Jika $t=\dfrac{3\pi}{4}$,
maka $(x+y)\sqrt{2}$ adalah \ldots
\begin{pilB}
  \item $25$ \item $35$ \item $45$ \item $50$ \item $60$
\end{pilB}
\end{soal}

%---------- 35: Isian ----------
\begin{soal}{35}
Misalkan $P$ adalah bidang yang melalui tiga titik $A(2,8,3)$,
$B(2,0,0)$, dan $C(0,2,0)$. Titik $X(2,2,w)$ berada pada bidang
$P$ jika nilai $w$ adalah \ldots
\end{soal}

%---------- 36: Isian ----------
\begin{soal}{36}
Diketahui vektor-vektor $\vec{a}=(1,x,1)$, $\vec{b}=(y,3,-5)$,
dan $\vec{c}=(xy,-12,z)$. Jika vektor $\vec{a}$ tegak lurus
dengan vektor $\vec{b}$ dan vektor $\vec{b}$ sejajar dengan
vektor $\vec{c}$, maka nilai $y+z$ adalah \ldots
\end{soal}

%---------- 37: Isian ----------
\begin{soal}{37}
Jika
\[
  \int_0^3 5f(x)\,dx+\int_0^1 8f(x)\,dx=18
  \qquad\text{dan}\qquad
  \int_0^3 5f(x)\,dx+\int_1^3 8f(x)\,dx=18,
\]
maka tentukan
\[
  \int_1^3 f(x)\,dx.
\]
\end{soal}

%---------- 38: PG Biasa ----------
\begin{soal}{38}
Jika $x_1$ dan $x_2$ adalah dua solusi persamaan
\[
  x^{2\log x-24}=\frac{x^{10}}{10^5},
\]
maka $x_1 x_2=10^k$ dengan $k$ adalah \ldots
\begin{pilB}
  \item $5$ \item $7$ \item $12$ \item $17$ \item $19$
\end{pilB}
\end{soal}

%---------- 39: PG Biasa ----------
\begin{soal}{39}
Misalkan $S_n$ adalah jumlah parsial $n$ suku pertama suatu deret
geometri. Jika $S_2=16$ dan $S_4=160$, maka suku ke-2 deret
tersebut adalah \ldots
\begin{pilB}
  \item $4$ \item $8$ \item $12$ \item $16$ \item $24$
\end{pilB}
\end{soal}

% ===================================================================
%  SOAL BARU  (nomor 40)
% ===================================================================

%---------- 40: PG Biasa ----------
\begin{soal}{40}
Diketahui fungsi $f(x)=2x^3-3x^2-12x+4$. Nilai maksimum lokal
dari fungsi tersebut adalah \ldots
\begin{pilB}
  \item $-16$ \item $-4$ \item $4$ \item $11$ \item $15$
\end{pilB}
\end{soal}

\end{document}
